(1) 每个动量态占据 \(k\) 空间体积 \((2\pi)^3/V\)，且具有两个自旋态。将半径 \(k_F\) 的球全部填满，粒子数 \(N=2[V/(2\pi)^3](4\pi k_F^3/3)=Vk_F^3/(3\pi^2)\)，故 \(k_F=(3\pi^2N/V)^{1/3}\)，具有逆长度单位。
(2) 单粒子能量为 \(\epsilon_k=\hbar^2k^2/(2m)\)，Fermi能 \(E_F=\hbar^2(3\pi^2n)^{2/3}/(2m)\)，其中 \(n=N/V\)。连续态求和 \(U=2[V/(2\pi)^3]4\pi\int_0^{k_F}\hbar^2k^4/(2m)\,dk=(3/5)NE_F\)，其平均动能为Fermi能的 \(3/5\)。
(3) 固定 \(N\) 时 \(E_F\propto V^{-2/3}\)，所以总能量 \(U=(3/5)NE_F\propto V^{-2/3}\)。由热力学定义 \(P=-(\partial U/\partial V)_N=(2/3)U/V=(2/5)nE_F\)，体现零温仍存在由Pauli排斥和动量填充造成的简并压力。
(4) 因为 \(P\propto n^{5/3}\propto V^{-5/3}\)，有 \(K=-V(\partial P/\partial V)_N=(5/3)P=(2/3)nE_F\)。高密度下压强按 \(n^{5/3}\) 增长，体模量同样按此幂律增加；在高密度相对论区需更换非相对论能量色散。
